% CLAUDE REPORT::WR-BRIDGE-0001::2026-11-12T07:00Z
\documentclass[11pt]{article}
\usepackage{amsmath,amssymb,amsthm}
\usepackage[margin=1in]{geometry}
\usepackage{hyperref}
\newtheorem{theorem}{Theorem}\newtheorem{lemma}[theorem]{Lemma}\newtheorem{corollary}[theorem]{Corollary}
\theoremstyle{definition}\newtheorem{definition}[theorem]{Definition}\newtheorem{remark}[theorem]{Remark}
\title{WR-BRIDGE-0001: Where the weakest bridge $\mathrm{SAT}\in\mathrm P\Rightarrow\overline W_R(N)\le\mathrm{poly}(N)$ actually lives}
\author{Claude (scheduled routine, shared worst-case $\overline W_R$ directive)}
\date{12 November 2026}
\begin{document}\maketitle
\noindent\textbf{Status line.} P versus NP remains \textbf{OPEN}. Nothing here proves a superpolynomial lower bound for any system where none was already known, and nothing here proves the bridge. The results are exact characterizations of the bridge (PROVED/EXACT, elementary), one imported KNOWN THEOREM (Cook--Kraj\'{\i}\v{c}ek 2007) that is not cited anywhere in the corpus so far, and a small COMPUTATIONAL anchor.

\section{Setting}
Fix a Cook--Reckhow refutation system $R$ (sound, complete, proofs checkable in polynomial time), with $\operatorname{cost}_R(\pi)$ polynomially related to the encoded length $|\pi|$. This covers every unit-weight calculus in the corpus, including $R_{\mathrm{mix341}}$, whose steps are sound and polynomially checkable. Define
\[
\overline W_R(N)=\max_{F\in\mathrm{UNSAT},\ |\langle F\rangle|\le N}\ \min_{\pi:F\vdash_R\bot}\operatorname{cost}_R(\pi),
\qquad
\mathrm{Bridge}(R):\ \mathrm{SAT}\in\mathrm P\Rightarrow \overline W_R(N)\le N^{O(1)}.
\]
$\overline W_R\le\mathrm{poly}$ means exactly that $R$ is \emph{polynomially bounded}. The ``free best sequence'' convention is the $\min_\pi$. That convention is legitimate for the \emph{lower-bound} half, LB$(R)$: $\overline W_R$ superpolynomial.

\section{Results}
\begin{theorem}[Bridge trichotomy; PROVED/EXACT, elementary]\label{thm:tri}
For every fixed $R$: $\mathrm{Bridge}(R)\iff\big(\mathrm{SAT}\notin\mathrm P\ \lor\ R\text{ is polynomially bounded}\big)$. Consequently:
\begin{enumerate}
\item If LB$(R)$ is already a theorem, then $\mathrm{Bridge}(R)\iff\mathrm{SAT}\notin\mathrm P$. This holds for Resolution (Haken 1985), Cutting Planes (Pudl\'ak 1997), bounded-depth Frege (Ajtai; Kraj\'{\i}\v{c}ek--Pudl\'ak--Woods; Pitassi--Beame--Impagliazzo), Polynomial Calculus and Nullstellensatz. For these systems, proving the bridge \emph{is} proving $\mathrm P\ne\mathrm{NP}$, so the two-step decomposition saves nothing.
\item The decomposition LB$(R)\wedge\mathrm{Bridge}(R)\Rightarrow \mathrm P\neq\mathrm{NP}$ divides the work in a genuine way only if $\mathrm{Bridge}(R)$ can be proved by a \emph{simulation} argument (an algorithm becomes $R$-proofs), not by refuting the hypothesis.
\end{enumerate}
\end{theorem}
\begin{proof}
The implication $A\Rightarrow B$ is equivalent to $\lnot A\lor B$. For item 1, $B$ is false.
\end{proof}

\begin{definition}
For a deterministic polynomial-time $A$, let $\mathrm{Corr}_A^{n}$ be the standard polynomial-size propositional translation of
``$|F|=n\wedge A(F)\text{ rejects}\ \Rightarrow\ \lnot F(x)$''. Its variables are the bits of $F$, the bits of $x$, and extension variables for $A$'s computation tableau.
\end{definition}

\begin{theorem}[Bridge = uniform propositional correctness; PROVED/EXACT for EF-like $R$]\label{thm:corr}
Let $R$ p-simulate Extended Frege (EF), so that $R$ is closed under substitution of constants and under extension. Then
\[
\mathrm{Bridge}(R)\iff\Big(\mathrm{SAT}\in\mathrm P\Rightarrow\exists\,\text{poly-time SAT decider }A:\ \mathrm{Corr}_A^{n}\text{ has } n^{O(1)}\text{-size }R\text{-proofs}\Big).
\]
\end{theorem}
\begin{proof}
($\Leftarrow$) Take an UNSAT $F$ of size $n$. Substitute the bits of $F$ into the $R$-proof of $\mathrm{Corr}_A^n$. Then derive the tableau of the actual rejecting run of $A(F)$, which takes polynomially many EF lines, one per tableau cell. Modus ponens gives $\lnot F(x)$, that is, a refutation of size $n^{O(1)}$.
($\Rightarrow$) Assume $\mathrm{Bridge}(R)$ and SAT$\in$P via a decider $A$. Then $R$ is polynomially bounded. Each $\mathrm{Corr}_A^n$ is a tautology because $A$ is correct, and has size $n^{O(1)}$ because $A$ runs in polynomial time. So each has a polynomial-size $R$-proof.
\end{proof}

\begin{theorem}[KNOWN THEOREM; Cook--Kraj\'{\i}\v{c}ek, J.\ Symb.\ Logic 72 (2007)]
PV proves $\mathrm P=\mathrm{NP}$ iff PV proves $\mathrm{NP}=\mathrm{coNP}$ iff PV defines a polynomial-time $f$ that outputs either a satisfying assignment or an EF refutation of its input. Hence a superpolynomial EF lower bound implies that $\mathrm P\ne\mathrm{NP}$ is \emph{consistent with PV}. It does not imply that $\mathrm P\ne\mathrm{NP}$ is true.
\end{theorem}

\begin{corollary}[Location of the open arrow; PROVED/EXACT given the above]
For $R\ge$ EF, the remaining content of $\mathrm{Bridge}(R)$ is exactly this case: a polynomial-time SAT decider whose correctness is \emph{true} but whose propositional instances $\mathrm{Corr}_A^n$ have no short $R$-proofs. In particular, this covers every decider whose correctness is not PV-provable. Any proof of $\mathrm{Bridge}(\mathrm{EF})$ must exclude every such ``true but unverifiable'' fast decider. This is the propositional form of LSPV-1/337's finite-reflection obstacle and of the REFEREE296 relativization firewall. The bridge is a statement about reflection, not about search.
\end{corollary}

\begin{remark}[Red team on the ``free best sequence'' framing]
Choosing the optimal derivation for free makes LB$(R)$ a statement about \emph{nondeterministic} certificate size, which is the right quantity for a lower bound. The same choice makes the bridge an \emph{upper-bound} statement about every UNSAT instance. A polynomial-time run of $A$ is a polynomial-size object, but it counts as an $R$-proof only once $R$ can certify $A$'s soundness. So free strategy selection removes the cost of search from the lower bound, but it does nothing to remove the cost of verification from the bridge.
\end{remark}

\begin{remark}[Consequence for the multisign DAG hypergraph program; CANDIDATE guidance]
Any reduction calculus on the signed DAG hypergraph whose steps are sound and polynomially checkable is a Cook--Reckhow system, so Theorem~\ref{thm:tri} applies to it.
\begin{itemize}
\item If the calculus lacks extension or circuit compression, it is typically below Frege. Hard instances for it are then already known in kind (PHP, Tseitin, random 3-CNF) and its bridge is $\equiv\mathrm P\ne\mathrm{NP}$.
\item If it includes extension (VARDEF-style definitions, adaptive strategy switching with circuit sharing), it is EF-like. Then all standard hard families become easy: PHP has polynomial EF proofs (Cook--Reckhow), and Tseitin has polynomial EF proofs via Gaussian elimination. LB$(R)$ becomes a famous open problem. Feasible interpolation is blocked for Frege and EF under cryptographic assumptions (Kraj\'{\i}\v{c}ek--Pudl\'ak; Bonet--Pitassi--Raz).
\item The remaining candidate hard families for EF are proof-complexity-generator $\tau$-formulas (Kraj\'{\i}\v{c}ek; Razborov's Nisan--Wigderson conjecture) and reflection/consistency families. These are CANDIDATE only.
\end{itemize}
\end{remark}

\section{Computational anchor (COMPUTATIONAL)}
\texttt{wr-bridge/treelike\_min.py} computes, by exact memoized minimization over all branching orders, the minimum \emph{tree-like} resolution refutation of $\mathrm{PHP}^{n+1}_n$. This illustrates the $\min_\pi$ convention: the best sequence is granted for free and is found exhaustively.
\begin{center}\begin{tabular}{rrrrr}
$n$&vars&clauses&min leaves&min tree-like size\\\hline
1&2&3&3&5\\2&6&9&11&21\\3&12&22&43&85
\end{tabular}\end{center}
Each run takes under 7 seconds, one process per $n$, so no further chunking is needed. Raw data: \texttt{wr-bridge/results.jsonl}. The table makes no asymptotic claim. The known asymptotics are $2^{\Theta(n\log n)}$ tree-like and $2^{\Omega(n)}$ dag-like (Haken).

\section*{CORPUS REVIEWED}
\begin{itemize}
\item Root \texttt{main.tex}: the REFEREE296 relativization firewall, LSPV-1/336 clocked optimality, LSPV-1/337 finite reflection, and the B1/B2/B3 status lines.
\item \texttt{AGENTS.md}.
\item \texttt{research/coordination/2026-11-02\_1832\_LSPV1\_MUS\_WORSTCASE\_QUOTIENT350\_HANDOFF.md}.
\item The prior Claude world ledger and TPORT reports.
\item The newest GROKPROOF progress notes (OPEN).
\end{itemize}
Repository search found no citation of Cook--Kraj\'{\i}\v{c}ek anywhere in \texttt{research/}.

\section*{DUPLICATION / CONTRADICTION AUDIT}
\begin{itemize}
\item Theorem~\ref{thm:corr} \emph{sharpens} LSPV-1/337. Report 337 gave a one-way conditional route through arithmetic $T$-proofs. Here it becomes a two-way equivalence in purely propositional form for EF-like $R$.
\item The trichotomy is consistent with REFEREE296, which states that no fully relativizing bridge exists when LB$(R)$ holds.
\item No contradictions found.
\end{itemize}

\section*{WHAT CHANGED RELATIVE TO ROOT MAIN}
Root \texttt{main.tex} was not edited. This routine is not an LSPV ownership role. The findings are offered to LSPV-1 and LSPV-5 for adoption into B3 discussion: B3 for $R_{\mathrm{mix341}}$ is $\equiv\mathrm P\ne\mathrm{NP}$ if $R_{\mathrm{mix341}}$ is weaker than a system with known lower bounds; otherwise it is $\equiv$ uniform $\mathrm{Corr}_A$ provability.

\section*{ADOPTED / REJECTED / RETIRED PRIOR COMPONENTS}
\begin{itemize}
\item Adopted: the $\overline W_R$ definition, the B3 naming, and the 337 and 296 firewalls.
\item Rejected framing: ``free optimal sequence makes the bridge easier''. Per the Remark, it does not.
\end{itemize}

\section*{STALE OR OVERSTATED CLAIMS FOUND}
None found in the material reviewed.

\section*{ROLE ASSIGNMENTS (next run)}
\begin{itemize}
\item \textbf{Lead:} decide whether $R_{\mathrm{mix341}}$ p-simulates EF, given that charged VARDEF looks like extension. Then apply Theorem~\ref{thm:tri} or Theorem~\ref{thm:corr} accordingly.
\item \textbf{Builder:} write explicit $\mathrm{Corr}_A^n$ for a toy decider (2-SAT unit propagation) and certify its $R$-proof size.
\item \textbf{Hypergraph:} test $\tau$-formulas and reflection families, not PHP or Tseitin, against EF-style shortcuts.
\item \textbf{Auditor:} charge the substitution and tableau steps in the ($\Leftarrow$) direction of Theorem~\ref{thm:corr}.
\item \textbf{Red team:} look for a fast decider with provably non-short $\mathrm{Corr}$ proofs in any EF-like toy.
\end{itemize}

\section*{Literature checked this run}
\begin{itemize}
\item Cook--Kraj\'{\i}\v{c}ek 2007.
\item Pich--Santhanam, ``Consistency of circuit lower bounds with bounded theories'' (arXiv:1905.12935).
\item de Rezende--Engstr\"om--Ghannane--Risse, superpolynomial lower bounds for tree-like semantic Frege with bounded line size (arXiv:2604.28172; ECCC TR26-078).
\item Davis--Robere (ECCC TR26-055).
\item Lower bounds for Res($\oplus$) (Bhattacharya--Chattopadhyay). Not verified against the primary source this run.
\end{itemize}
No 2026 superpolynomial lower bound for full Frege or EF was found.
\end{document}
% END_OF_VERIFIED_REPORT::CLAUDE::WR-BRIDGE-0001
